\chapter{Discussion}
\label{ch:discussion}

The results support a deliberately bounded claim. This chapter separates what the
framework demonstrates strongly from what it demonstrates only partially, and argues
that the distinction turns on two different questions that prior work often conflates.
The central thesis is this: the framework contributes a \emph{validated demand-gap
diagnostic} and a \emph{characterized, partially-successful corridor generator}, and
these two contributions carry different evidential weight.

\section{Two questions, two layers}
\label{sec:disc-twoquestions}

It is useful to distinguish two questions that a data-driven planning framework might
be asked to answer. Question~A: does the generator \emph{reproduce corridors that were
actually built}? Question~B: does the generator \emph{produce good corridors}---lines
that serve real, unmet demand efficiently and without redundancy? Much of the
literature on machine-learning station siting implicitly answers a version of
Question~A, treating agreement with the existing network as validation. But A is a
weak and asymmetric proxy for B, because built lines are not guaranteed to be optimal:
they are chosen under political, financial, and land-availability constraints that the
framework does not model. Agreement therefore corroborates; disagreement is faint
evidence of error. The two pipeline layers map onto these questions---the diagnostic
layer (W1--W4) is validated against B-type evidence, while the generator (W5--W6) is
tested against both---and they succeed to different degrees.

\section{The diagnostic layer is the strong contribution}
\label{sec:disc-diagnostic}

The coverage-gap diagnosis (\cref{ch:methodology}, \cref{sec:w3}) is the framework's
firmest result. Its methodological contribution is to break the circularity that
undermines ``has-a-stop'' targets: by defining the dependent variable as the ratio of
independently-estimated demand to \textsc{gtfs}-measured supply
(\cref{eq:gap}), the model learns where transit is \emph{needed} rather than where it
\emph{exists}. The supervised retrain confirms this target is learnable from
place characteristics alone, with no supply-side leakage (\cref{tab:retrain}).

Three lines of evidence support the diagnostic. First, and most importantly, it is
corroborated \emph{out of sample}: the 2024 \textsc{gtfs} snapshot predates SITEUR
Line~4, so the corridor is treated as unserved, yet its \ageb{}s show \num{1.6} times
the metropolitan High-gap rate (\cref{sec:res-w8}). A corridor that professional
planners independently chose to build was flagged as high-priority by the model that
never saw it. Second, the diagnostic \emph{transfers}: it produces a plausible,
differentiated reading of unmet demand in two further cities (\cref{ch:transferability}),
ordering them sensibly by High-gap share. Third, it is \emph{interpretable}: SHAP
attribution identifies dense, high-need, employment-rich but under-served zones,
matching the geographic pattern in \cref{fig:coverage-gap}. This layer is the
defensible, citable core of the thesis.

\section{The generative layer is a characterized, partial success}
\label{sec:disc-generative}

The corridor generator is a weaker but genuine contribution, and its history is
instructive. In its first form the generator was essentially negative: of its feasible
corridors, only a degenerate short stub passed the merit tests, and feasibility was
confounded with anchor-cluster sparsity---the feasibility filter selected for
\emph{few, near-collinear anchors} rather than for good corridors. This was traced to
two coupled defects: a shaper that flattened a minimum spanning tree into a line with
phantom jumps, and a detour metric measured against the endpoint distance, which
over-penalizes a demand-coverage corridor that legitimately curves.

Re-architecting the generator around frontier anchors, a minimum-spanning-tree
\emph{diameter} trunk stitched from real roads, and an \emph{anchor-directness}
feasibility gate (\cref{sec:w6}) changed the verdict. On its own terms the generator
now produces at least one substantive, feasible, merit-passing corridor:
\code{W6\_G02} serves \SI{12.1}{\kilo\metre} and \num{25} \ageb{}s, of which
\SI{56}{\percent} are High-gap (versus the \SI{20.7}{\percent} baseline), is
non-redundant with the existing network (best Jaccard \num{0.18}), and sits at the
\num{73}rd percentile of demand per kilometre. This is a real line, not a stub.

The success is partial, however, and the thesis does not overstate it. Alongside
\code{W6\_G02}, the generator also surfaces lower-efficiency connectors
(\code{W6\_G00} and \code{W6\_G01}) that serve real need but rank poorly on demand per
kilometre. The residual limitation is architectural---the fixed anchor count, spatial
clustering, and spanning-tree trunk cannot target every high-gap pocket---and it is
now measured rather than merely asserted.

\section{Reconstruction is a weak proxy for corridor merit}
\label{sec:disc-reconstruction}

The clearest illustration of the Question-A/Question-B distinction is the generator's
failure to reconstruct built rail. Across masked backtests it does not re-discover the
premium lines: masking the whole premium tier yields \SI{15}{\percent} mean shape
overlap, and masking individual rail lines yields essentially zero
(\cref{sec:res-w8}); the Line~4 corridor is matched at only \num{0.05} recall. Taken
naively, this looks like failure. Read correctly, it is close to the expected outcome
and only weak evidence against the generator.

The reason is mechanical as well as epistemic. Masking a single rail line barely moves
accessibility, because parallel bus routes remain and the served/unserved seam scarcely
shifts; no strong new gap forms on the held-out corridor, so the frontier generator has
nothing to anchor on. This is the same mechanism that makes the demand-trunk backtest
degenerate in bus-only Toluca (\cref{sec:transf-validation}). Epistemically, even a
perfect reconstruction score would not establish corridor merit, because the built
lines encode non-demand constraints. The benchmark in the transfer cities is the
instructive mirror image: where the existing network already covers demand well, the
generator overlaps it heavily (Toluca \SI{75}{\percent}, Aguascalientes
\SI{54}{\percent})---revealed-preference agreement, not new coverage. The honest
framing, therefore, is that the generator's value lies in Question~B, where it is
partially validated, not in Question~A, where agreement or disagreement is only weakly
diagnostic.

\section{Equity implications}
\label{sec:disc-equity}

Equity is treated as an explicit, separable term rather than an implicit consequence.
The additive formulation (\cref{eq:final}) lets the thesis report prioritization with
and without the equity bonus, and the sensitivity analysis shows the ranking is robust
to the equity weight ($\alpha$): the broad structure is fixed, and $\alpha$ shifts only
which specific \ageb{}s top the list. Substantively, the diagnostic and the equity term
point the same way---High-gap areas are disproportionately dense, high-social-lag
peripheries---so prioritizing unmet demand and prioritizing disadvantage are, in \zmg,
largely aligned rather than in tension. The correction of the marginalization index's
sign and imputation (documented in the data chapter) matters here: an equity term that
had silently rewarded the least-marginalized areas would have inverted this conclusion.

\section{Limitations}
\label{sec:disc-limitations}

Several limitations bound the results beyond those already noted. The demand surface
rests on a transferred distance-decay parameter and Euclidean impedance; a
network-distance gravity model calibrated on local survey data would tighten absolute
magnitudes. The \textsc{gtfs} snapshot is a single point in time, so the supply layer
does not capture service-frequency variation through the day. The generator's
feasibility gate is sensitive to source-\textsc{gtfs} stop density, which in the audit
manifests as an artifactually low feasible count for finely-spaced concessioned
networks. Most importantly, the \emph{generated} corridors are validated only
indirectly---through need, non-redundancy, and demand efficiency---and not against
observed ridership, because no counterfactual ridership exists for lines that were
never built. The diagnostic layer, by contrast, has an out-of-sample anchor in Line~4.
This asymmetry is the honest boundary of the contribution: a validated diagnostic, and
a characterized generator whose effectiveness is measured on its own terms.
